%% kapa: exam=1 round=21 year=2010 period=1 form=A subject=civil_law subject_ko=민법 seq=9 number=9
%% source: Q-Net 감정평가사 기출문제 (cst00309) artlSeq=5024181 file=extracted/1008/5024181/2010년 제21회 감정평가사 1차시험 기출문제/1교시.hwp
%% conversion: kordoc 4.24.51 -> tools/gen1.py (auto); review=proofread-sonnet; flags=none
\begin{question}{9}
다음 중 유효한 법률행위를 모두 고른 것은? (다툼이 있으면 판례에 의함)
\begin{framedtext}
ㄱ. 허가 없이 음식점을 경영하는 사람과 체결한 음식물매매계약 \par
ㄴ. 제2매수인이 매도인의 배임행위에 적극 가담하여 체결한 부동산 이중매매계약 \par
ㄷ. 토지거래허가 구역 내의 토지에 대한 매매계약의 정지조건이 허가 전에 불성취로 확정된 경우, 그 매매계약 \par
ㄹ. 부첩관계를 단절하면서 그 자녀의 생활비와 교육비를 지급하기로 하는 계약 \par
ㅁ. 세금을 감액 받을 목적으로 매매를 가장하여 체결한 증여계약 \par
\end{framedtext}
\begin{choices}
  \item ㄱ, ㄹ
  \item ㄴ, ㄷ
  \item ㄷ, ㅁ
  \item ㄱ, ㄹ, ㅁ
  \item ㄴ, ㄹ, ㅁ
\end{choices}
\end{question}
\answer{4}
